\documentclass[11pt,letterpaper]{article}

\usepackage{graphicx}           % Graphics package
\usepackage{amsmath} 
\usepackage{mathrsfs}
\usepackage[colorlinks=true]{hyperref} 
\usepackage{color}
\usepackage{amsfonts}
\usepackage[textheight=9in, textwidth=7.5in, letterpaper]{geometry}
\usepackage[makeroom]{cancel}
\usepackage{cite}
\usepackage{sectsty}
\usepackage{empheq}
\usepackage{appendix}
\usepackage{enumerate} 
\usepackage{simpler-wick}
\usepackage{amssymb}

\sectionfont{\fontsize{12}{12}\selectfont}
\subsectionfont{\fontsize{12}{12}\selectfont}

\newcommand*\widefbox[1]{\fbox{\hspace{2em}#1\hspace{2em}}}
\newcommand{\LP}{\left(}
\newcommand{\RP}{\right)}
\newcommand{\mc}[1]{\mathcal{#1}}
\newcommand{\R}{\mathbb{R}}
\newcommand{\C}{\mathbb{C}}
\newcommand{\Z}{\mathbb{Z}}
\newcommand{\D}{\partial}
\newcommand{\KET}[1]{\left| #1 \right\rangle }
\newcommand{\BRA}[1]{\left\langle #1 \right| }
\newcommand{\IP}[2]{\left\langle #1 \middle| #2 \right\rangle}
\newcommand{\PA}[2]{\frac{\partial #1}{\partial #2}}



%%%%%%%%%%%%%%%%---------------------MOST USEFUL COMMAND EVER---------------------%%%%%%%%%%%%%%%%%%%%%%
\newcommand{\fixme}[1]{{\bf {\color{red}[#1]}}}
\newcommand{\draftmode}{\usepackage[notref,notcite]{showkeys}}
\providecommand*\showkeyslabelformat[1]{\normalfont\sffamily\footnotesize#1}

%%%%%%%%%%%%%%%


\setcounter{tocdepth}{4}
\setcounter{secnumdepth}{4}

\begin{document}

\title{Notes on Large Gauge Transformations}

\author{Mathew Calkins}

\date{\today}

 
\maketitle

\abstract{Notes on large gauge transformations}

\tableofcontents

\section{Electrodynamics}

We work in the conventions of Daniel Harlow's excellent notes on QFT. In particular, we work in mostly-plus signature.

\subsection{Set Up}
The Maxwell Lagrangian density \begin{equation}\label{maxwelllagrangian}
\mathcal{L} = - \frac{1}{4} F_{\mu \nu} F^{\mu \nu} + A_\mu J^\mu
\end{equation}
gives equations of motion \begin{equation}
\D_\nu F^{\mu \nu} = J^\mu
\end{equation}
and $F = dA$ for free satisfies the Bianchi identity \begin{equation}
\D_\alpha F_{\beta \gamma} + \mbox{(cyclic permutations)} = 0.
\end{equation}
We denote \begin{equation}
E_i = F_{i0}, \quad B_i = \frac{1}{2} \epsilon_{ijk} F_{jk}.
\end{equation}
Under a gauge transformation \begin{equation}
A \to A' \equiv A + d \Omega, \quad A_\mu \to A_\mu + \D_\mu \Omega
\end{equation}
we have immediately $F \to F$, so the Lagrangian density shifts by \begin{equation}
A'_\mu J^\mu - A_\mu J^\mu = (\D_\mu \Omega) J^\mu = \D_\mu \LP \Omega J^\mu \RP - \Omega \D_\mu J^\mu.
\end{equation}
So in order for our theory to be consistent, the background source $J$ must obey \begin{equation}
\D_\mu J^\mu = 0.
\end{equation}







\subsection{Canonical Picture}
We make a few observations about the canonical structure. First, $\dot{A}_0$ appears nowhere in $\mathcal{L}$, so that \begin{equation}
\Pi^0(x) = \PA{\mc{L}}{\dot{A}_0}(x) = 0.
\end{equation}
This vanishing isn't the result of an equation of motion; the definition of $\Pi^0(x)$ sets it to 0. This is inconsistent with the canonical relation \begin{equation}
[A_0(\vec{x}), \Pi^0(\vec{y})] \stackrel{?}{=} i \delta^{(3)}(\vec{x} - \vec{y}).
\end{equation}
The spatial parts work more nicely. We find \begin{equation}
\Pi^i(x) = \PA{\mc{L}}{\dot{A}_i}(x) = F^{i0}(x) = - E^i(x).
\end{equation}
Then our Lagrangian equations of motion in terms of the canonical variables and their momenta have $\mu = 0$ component \begin{equation}
J^0 = \D_\nu F^{0 \nu} = \D_i F^{0i} \to \rho = \D_i E^i = - \D_i \Pi^i.
\end{equation}
So we have a non-dynamical Gauss law constraint. This does not determine the time-evolution, but constrains the physically allowed configurations at any fixed time. We can think of the law $\Pi^0 = 0$ similarly. \begin{equation}
\Pi^0 = 0, \quad \D_i \Pi^i + \rho = 0.
\end{equation}
It may not be clear why these warrant being treated on equal footing, when the first came from the definition of momentum while the second came from the equations of motion. Regardless, we shall interpret these as constraints on the phase space: time-evolution will be governed by Hamilton's equations of motion in the usual way, and the constraints cut phase space down to a physical subspace of legal instantaneous states. When we compute the transformations generated by various charges in the usual way via Poisson brackets, we will only get sensible results if we work out all brackets before (if at all) reducing modulo these constraints. If this is all consistent, then time evolution generated by Poisson brackets against the Hamiltonian should take legal states to legal states.\\
\\
\
Under the canonical relations \begin{equation}
[A_\mu(\vec{x}), \Pi^\nu(\vec{y})] = i \delta_\mu^\nu  \delta^{(3)} (\vec{x} - \vec{y})
\end{equation}
these quantities that vanish under our constraints also turn out to generate gauge transformations. Concretely, we compute \begin{equation}
\begin{aligned}
\left[ \int d^3 \vec{y} \ \dot{\Omega}(\vec{y}) \Pi^0(\vec{y}), A_0 (\vec{x}) \right] & = \int d^3 \vec{y} \ \dot{\Omega}(\vec{y}) \left[  \Pi^0(\vec{y}), A_0 (\vec{x}) \right] \\
& = - i \int d^3 \vec{y} \ \dot{\Omega}(\vec{y}) \delta^{(3)}(\vec{x} - \vec{y}) \\
& = - i \dot{\Omega}(\vec{x}).
\end{aligned}
\end{equation}
Similarly, \begin{equation}
\begin{aligned}
\left[ \int d^3 \vec{y} \ \Omega(\vec{y}) \LP \D_j \Pi^j(\vec{y}) + \rho(\vec{y}) \RP, A_i (\vec{x}) \right] & = \int d^3 \vec{y} \ \Omega(\vec{y}) \left[  \D_j \Pi^j(\vec{y}), A_i (\vec{x}) \right] \\
& = \int d^3 \vec{y} \ \D_{y^j} \LP \Omega(\vec{y}) \left[ \Pi^j(\vec{y}), A_i (\vec{x}) \right] \RP - \int d^3 \vec{y} \ \LP \D_j \Omega(\vec{y}) \RP \left[   \Pi^j(\vec{y}), A_i (\vec{x}) \right] \\
& = -i \int d^3 \vec{y} \ \D_{y^j} \LP \Omega(\vec{y}) \delta^j_i \delta^{(3)}(\vec{x} - \vec{y}) \RP + i \int d^3 \vec{y} \ \LP \D_j \Omega(\vec{y}) \RP \delta^i_j \delta^{(3)}(\vec{x} - \vec{y}) \\
& = i \D_i \Omega(\vec{x})
\end{aligned}
\end{equation}
where in the last line we have assumed that $\Omega(\vec{y})$ goes to 0 at spatial infinity. Soon we will remove this assumption and see what we get.


\subsection{Quantization}
In the classical theory we said that the legal states were the points in phase space satisfying $\Pi^0 = \D_i \Pi^i + \rho = 0$. The quantum-mechanical analogue is that states should be annihilated by the corresponding operators. This is perhaps easiest to understand if we quantize by first building an over-large Hilbert space of functionals $\Psi[A] \in \mc{H}_{\mbox{\tiny{big}}}$. In this realization of our Hilbert space, our canonical momenta act as \begin{equation}
\Pi^\mu(\vec{x}) = -i \frac{\delta}{\delta A_\mu(\vec{x})}.
\end{equation}
The states annihilated by $\Pi^0$ are exactly those functionals which happen to be independent of $A_0$. The remaining constraint then becomes \begin{equation}
\rho(\vec{x}) \Psi[\vec{A}] = - \D_{x^j} \Pi^j(\vec{x}) \Psi[\vec{A}] = i \D_{x^j} \frac{\delta}{\delta A_j(\vec{x})} \Psi[\vec{A} ].
\end{equation}
Then we can implement gauge transformations \begin{equation}
A \to A + d \Omega \implies A_i \to A_i + \D_i \Omega
\end{equation}
concretely on physical states as \begin{equation}
\begin{aligned}
\Psi[\vec{A} + \nabla \Omega] & = \exp \left\{ \int d^3 \vec{x} \LP \D_j \Omega(\vec{x}) \RP \frac{\delta}{\delta A_j(\vec{x})} \right\} \Psi[\vec{A}] \\
& = \exp \left\{ i \int d^3 \vec{x} \LP \D_j \Omega(\vec{x}) \RP \Pi^j(\vec{x}) \right\} \Psi[\vec{A}] \\
& = \exp \left\{ - i \int d^3 \vec{x}  \Omega(\vec{x}) \LP   \D_j \Pi^j(\vec{x}) \RP \right\} \Psi[\vec{A}] \\
& = \exp \left\{ i \int d^3 \vec{x}  \Omega(\vec{x}) \rho(\vec{x})  \right\} \Psi[\vec{A}].
\end{aligned}
\end{equation}
Note that, at an intermediate step, we have integrated by parts on the fixed-time slice in question and assumed $\Omega$ (or more specifically, $\Omega \rho$ vanishes sufficiently quickly at spatial infinity. \\
\\
\
There is an interesting story we can tell here. First, we see that physically-admissible states $\Psi[\vec{A}]$ change under small gauge transformations, by only by a phase independent of the state in question. So the corresponding rays in Hilbert space, which are what matter at the end of the day, are unchanged. Even at the level of vectors, it is interesting to note that this phase is an artifact of our working in a theory where $A$ is dynamical but the background source $J$ is frozen. If we performed the same analysis in a theory sourced by dynamical charged matter, the corresponding degrees of freedom in the quantum theory would contribute the opposite phase, and the corresponding legal states would be exactly gauge-invariant. 



\subsection{Large Gauge Transformations}
Under a general variation $A \to A + \delta A$, the Lagrangian density \eqref{maxwelllagrangian} varies as \begin{equation}
\begin{aligned}
\delta \mc{L} & = \delta \LP - \frac{1}{4} F_{\mu \nu} F^{\mu \nu} + A_\mu J^\mu \RP \\
& = - F^{\mu \nu} \D_\mu \delta A_\nu + J^\mu \delta A_\mu \\
& = \LP \D_\mu F^{\mu \nu} + J^\nu \RP \delta A_\nu - \D_\mu \LP F^{\mu \nu} \delta A_\nu \RP
\end{aligned}
\end{equation}
We temporarily IR-regulated by placing spatial infinity $\Gamma$ at some finite point (say, $|\vec{x}| = R$). Then the variation of the action has a boundary term \begin{equation}
\delta S \supset - \int_\Gamma d^3 x n_\mu F^{\mu \nu} \delta A_\nu
\end{equation}
where $n^\mu$ is the outward-oriented normal vector to $\Gamma$. For the action to be stationary (setting aside contributions from timelike boundary components) we need \begin{equation}
n_\mu F^{\mu \nu} \delta A_\nu |_\Gamma = 0.
\end{equation}
We can accomplish this by restricting our domain of variation to $A$ such that \begin{equation}
n_\mu F^{\mu \nu} |_\Gamma = 0 \quad \mbox{(Neumann boundary conditions)}
\end{equation}
or equivalently such that the normal component of $\vec{E}$ and the tangential component of $\vec{B}$ vanish. Alternatively, if we work in terms of differential forms we find (up to constants depending on our normalization of the Hodge star) \begin{equation}
\begin{aligned}
\mc{L} d \mbox{Vol} & = - \frac{1}{2} F \wedge * F + A \wedge * J \\
\implies \delta \mc{L} d \mbox{Vol}  & = - (\delta d A) \wedge * F + \delta A \wedge * J \\
& = \delta A \wedge \LP d * F + * J \RP - d \LP \delta A \wedge * F \RP
\end{aligned}
\end{equation}
and our boundary contribution is \begin{equation}
\delta \mc{L} d \mbox{Vol} \sim \int_\Gamma (\delta A) \wedge * F
\end{equation}
so it suffices to demand that $\delta A$ vanish when pulled back to the submanifold $\Gamma$, \textit{i.e.}, that \begin{equation}
t^\mu A_\mu |_\Gamma = 0 \mbox{ for all } t^\mu \mbox{ tangent to } \Gamma \quad \mbox{(Dirichlet boundary conditions)}
\end{equation}
which is equivalent to killing the normal component of $\vec{B}$ and the tangential component of $\vec{E}$. The latter is more natural in QFT, where we typically take $A_\mu$ to vanish at infinity. In order for gauge transformations $A \to A + d \Omega$ to respect these boundary conditions, we must have \begin{equation}
\Omega|_\Gamma = \mbox{constant}
\end{equation}


\bibliography{bibfile}
\bibliographystyle{utphys}

\end{document}





%We saw above that, either in the classical theory or in the Heisenberg picture of the quantum theory, the first-class quantities $\Pi^0$ and $\D_i \Pi^i - \rho$ generated gauge transformations of observables. From this we see that the  In the quantum theory they should equivalently act on states to generate gauge transformations
